\section{Arsitektur Web dan HTTP}
Pemrograman web berjalan di atas arsitektur client-server, di mana browser bertindak sebagai client dan server menyediakan data serta layanan. Komunikasi utama terjadi melalui protokol HTTP yang mengatur request dan response. Pemahaman dasar tentang HTTP membantu Anda menelusuri mengapa halaman bisa ditampilkan atau gagal dimuat \cite{mdnLearn}.

HTTP memiliki metode yang berbeda untuk tujuan berbeda, seperti mengambil data, mengirim formulir, atau memperbarui sumber daya. Dalam praktik, Anda sering menggunakan GET untuk membaca data dan POST untuk mengirim data. Konsep status code juga penting untuk memahami hasil dari sebuah request.

\begin{table}[ht]
\centering
\begin{tabular}{|P{3cm}|P{9cm}|}
\hline
\textbf{Metode} & \textbf{Kegunaan Utama} \\
\hline
GET & Mengambil data atau halaman dari server \\
\hline
POST & Mengirim data, misalnya formulir pendaftaran \\
\hline
PUT & Memperbarui data yang sudah ada \\
\hline
DELETE & Menghapus data tertentu \\
\hline
\end{tabular}
\caption{Metode HTTP yang sering digunakan}
\end{table}

